% Part: incompleteness
% Chapter: introduction
% Section: historical-background

\documentclass[../../../include/open-logic-section]{subfiles}

\begin{document}

\olfileid{inc}{int}{bgr}

\olsection{ऐतिहासिक पार्श्वभूमी}

या विभागात अपूर्णतेच्या प्रमेयांची पार्श्वभूमी समजण्यासाठी उपयुक्त अशा
ऐतिहासिक घडामोडींची थोडक्यात चर्चा करू. विशेषतः, गणितीय तर्कशास्त्राच्या
इतिहासाचा अगदी स्थूल आढावा घेऊ आणि मग गणिताच्या पायाभरणीच्या इतिहासाविषयी
थोडे बोलू.

\begin{digress}
  ``गणितीय तर्कशास्त्र'' या संज्ञेला दोन अर्थ आहेत. ``गणितीय'' हा शब्द
  अभ्यासविषय दर्शवतो असे मानल्यास, ``गणिताचे तर्कशास्त्र'' म्हणजे गणितीय
  युक्तिवादाची तत्त्वे. तो अभ्यासाची पद्धत दर्शवतो असे मानल्यास,
  ``तर्कशास्त्राचे गणित'' म्हणजे युक्तिवादाच्या तत्त्वांचा गणितीय अभ्यास.
  पुढील विवेचनात गणितीय तर्कशास्त्र हे दोन्ही अर्थांनी, अनेकदा एकाच वेळी,
  येते.
\end{digress}

तर्कशास्त्राचा अभ्यास मूलतः ॲरिस्टॉटलपासून सुरू झाला. तो सुमारे
384--322 \textsc{bce} या काळात जगला. त्याच्या \emph{कॅटेगरीज}, \emph{प्रायर
  ॲनॅलिटिक्स} आणि \emph{पोस्टीरियर ॲनॅलिटिक्स} या ग्रंथांत वैज्ञानिक युक्तिवादाच्या
तत्त्वांचा पद्धतशीर अभ्यास आहे; त्यात संवाक्याचाही सखोल व
पद्धतशीर अभ्यास समाविष्ट आहे.

मध्ययुगभर विद्वत् तत्त्वज्ञानावर ॲरिस्टॉटलच्या तर्कशास्त्राचे वर्चस्व होते.
अठराव्या शतकातसुद्धा कांटने त्याचे तर्कशास्त्र परिपूर्ण असून त्यात सुधारणा
करण्याची गरज नाही, असे म्हटले. पण न्यायवाक्यरचनेचा सिद्धांत गणितीय
युक्तिवादाच्या अगदी वरवरच्या पैलूंपलीकडे काही दाखवण्यासाठी अत्यंत मर्यादित
आहे. त्याच्या एक शतक आधी, न्यूटनचा समकालीन लाइब्निझ याने तार्किक
युक्तिवादासाठी पूर्ण ``कलन'' असावे अशी कल्पना केली होती. त्याने अशा कलनाची
रचना करण्याच्या दिशेने काही प्राथमिक पावलेही टाकली; त्यातून मूलतः विधानीय
तर्कशास्त्राची एक आवृत्ती वर्णिली गेली.

एकोणिसावे शतक तर्कशास्त्रासाठी निर्णायक ठरले. 1854 मध्ये जॉर्ज बूलने
\emph{विचाराचे नियम} लिहिला; त्यातील विधानीय तर्कशास्त्राचा सखोल बीजगणितीय
अभ्यास आधुनिक मांडण्यांपासून फारसा दूर नाही. 1879 मध्ये गॉटलोब फ्रेगेने
\emph{Begriffsschrift} (संकल्पनालेखन) प्रसिद्ध केला. त्यात विधानीय
तर्कशास्त्राला संख्यापक आणि संबंध यांची भर घातली असल्याने प्रथम-क्रम
तर्कशास्त्र समाविष्ट होते. प्रत्यक्षात, फ्रेगेच्या तार्किक पद्धतींत
उच्च-क्रम तर्कशास्त्र आणि त्याहून अधिक गोष्टीही होत्या. \emph{अंकगणिताचे
  मूलभूत नियम} या ग्रंथात त्याने सर्व अंकगणित केवळ तार्किक गृहीतकांपासून
आपल्या Begriffsschrift मध्ये निष्पन्न करता येईल, हे दाखवण्याचा प्रयत्न
केला. दुर्दैवाने, रसेलने 1902 मध्ये दाखवून दिल्याप्रमाणे ही गृहीतके
विसंगत निघाली. तरी तो विसंगत स्वयंसिद्धक बाजूला ठेवल्यास, फ्रेगेने आधुनिक
तर्कशास्त्र जवळजवळ एकट्याने निर्माण केले; ही विलक्षण कामगिरी होती.
बूलनंतर पिअर्स आणि श्रॉडर यांच्यासह बीजगणितीय दृष्टिकोन असलेल्या इतर
विचारवंतांनीही संख्यापकाधारित तर्कशास्त्र स्वतंत्रपणे विकसित केले.

आता गणिताच्या पायाभरणीतील घडामोडींकडे वळू. तर्कशास्त्राची गणितात महत्त्वाची
भूमिका असल्यामुळे नुकत्याच वर्णिलेल्या घडामोडींशी त्यांचा मोठा संबंध आहे.
उदाहरणार्थ, सर्व गणित केवळ आपल्या तार्किक चौकटीवर आधारता येईल हे
दाखवण्याच्या स्पष्ट उद्देशाने फ्रेगेने आपले तर्कशास्त्र विकसित केले.
विशेषतः, कांटच्या मते गणितातील सत्ये अनुभवपूर्व \emph{संश्लेषक} होती;
फ्रेगेला त्याऐवजी ती अनुभवपूर्व \emph{विश्लेषक} आहेत, हे दाखवायचे होते.

अनेकांच्या मते गणिताचा खरा उदय ग्रीकांबरोबर झाला. सुमारे 300 इ.स.पू.
मध्ये लिहिलेला युक्लिडचा \emph{मूलतत्त्वे} हा ग्रंथ काटेकोरपणा आणि
अचूकता यांवर भर देणाऱ्या ग्रीक गणिताचे विकसित उदाहरण आहे. युक्लिडच्या
\emph{मूलतत्त्वे}मधील व्याख्या आणि सिद्धता आजही माध्यमिक शाळेच्या
भूमितीच्या पाठ्यपुस्तकांत जवळजवळ तशाच टिकून आहेत (जिथे अजून माध्यमिक
शाळेत भूमिती शिकवली जाते तिथे). गणितातच नव्हे तर तत्त्वज्ञानाच्या
शाखांतही काटेकोर युक्तिवादाचा आदर्श म्हणून गणितीय युक्तिवादाचे हे रूप
मानले गेले आहे. (स्पिनोझाने नैतिक आणि धार्मिक युक्तिवादसुद्धा
युक्लिडच्या शैलीत मांडले; ते पाहणे काहीसे विचित्र वाटते!)

न्यूटन आणि लाइब्निझ यांनी सतराव्या शतकात कलनाचा शोध लावला. (प्राधान्यावरून
शतकानुशतके तीव्र वाद चालला; पण आज बहुतेक अभ्यासकांच्या मते त्यांची
बहुतांश प्रगती स्वतंत्रपणे झाली.) कलनात, उदाहरणार्थ, अनंतसूक्ष्म
राशींच्या अनंत बेरजांबद्दल युक्तिवाद करावा लागतो. या वैशिष्ट्यांमुळे
बिशप बर्कलीने टीका केली: ईश्वरावरील विश्वास आपल्या काळातील गणितापेक्षा
कमी तर्कसंगत नाही, असे त्याचे म्हणणे होते. अठराव्या शतकात कलनाच्या
पद्धती मोठ्या प्रमाणावर वापरल्या गेल्या; उदाहरणार्थ, लिओनहार्ड
ऑयलरने अनंत बेरजांचा समावेश असलेली गणिते करून प्रभावी निष्कर्ष मिळवले.

एकोणिसाव्या शतकात गणितज्ञांनी कलनाला अधिक भक्कम पाया देऊन बर्कलीच्या
टीकेला उत्तर देण्याचा प्रयत्न केला. कॉशी, वायेरश्ट्रास, बोल्झानो आणि
इतरांच्या प्रयत्नांतून ``एप्सिलॉन आणि डेल्टा'' यांच्या साहाय्याने सीमा,
सातत्य, अवकलन आणि समाकलन यांच्या आजच्या व्याख्या निर्माण झाल्या; म्हणजेच
त्यात अनंतसूक्ष्म राशींचा संदर्भ उरला नाही. पुढे त्याच शतकात गणितज्ञांनी
वास्तव संख्यांसह कलनाचे सर्व पैलू नैसर्गिक संख्यांच्या आधारे स्पष्ट
करण्याचा अधिक प्रयत्न केला. (क्रोनेकर: ``पूर्ण संख्या ईश्वराने निर्माण
केल्या; उरलेले सर्व मानवाचे कार्य आहे.'') 1872 मध्ये डेडेकिंडने
``सातत्य आणि अपरिमेय संख्या'' लिहून वास्तव संख्या परिमेय संख्यांच्या
संचांच्या रूपात कशा ``रचता'' येतात हे दाखवले. (परिमेय संख्यांकडे, तुम्हाला
माहीत आहेच, नैसर्गिक संख्यांच्या जोड्यांच्या रूपात पाहता येते.) 1888 मध्ये
त्याने ``Was sind und was sollen die Zahlen'' (साधारणतः ``नैसर्गिक
संख्या म्हणजे काय आणि त्यांचे स्वरूप कसे असावे?'') लिहून नैसर्गिक संख्या
निखळ ``तार्किक'' संज्ञांत स्पष्ट करण्याचा प्रयत्न केला. 1887 मध्ये
क्रोनेकरने ``\"Uber den Zahlbegriff'' (``संख्येच्या संकल्पनेविषयी'')
लिहून सर्व गणितीय वस्तू पूर्ण संख्यांच्या आधारे दर्शवण्याविषयी सांगितले.
1889 मध्ये ज्यूझेप्पे पेआनोने नैसर्गिक संख्यांसाठी आकारिक, प्रतीकात्मक
स्वयंसिद्धके दिली.

एकोणिसाव्या शतकाच्या अखेरीस अनंताशी व्यवहार करण्याचे नवे धाडसही दिसले.
त्याआधी अनंत वस्तू आणि रचना (उदा., नैसर्गिक संख्यांचा संच) अतिशय
सावधपणे हाताळल्या जात. ``अनंत अनेक'' म्हणजे ``तुम्हाला हव्या तितक्या''
आणि ``सीमेत जवळ जाणे'' म्हणजे ``तुम्हाला हवे तितके जवळ जाणे'' असे समजले
जाई. पण गेऑर्ग कांटोरने अनंताचा प्रत्यक्ष अर्थाने विचार करणे शक्य आहे हे
दाखवले. कांटोर, डेडेकिंड आणि इतरांच्या कार्यामुळे गणिताची आज व्यापकपणे
स्वीकारली जाणारी सर्वसाधारण संचसिद्धान्तीय समज विकसित होऊ लागली.

यातून आपण तर्कशास्त्र आणि पायाभरणीतील विसाव्या शतकाच्या घडामोडींपर्यंत
पोहोचतो. 1902 मध्ये रसेलला फ्रेगेच्या तार्किक पद्धतीत विरोधाभास सापडला.
1904 मध्ये झर्मेलोने तथाकथित ``निवडीचे स्वयंसिद्धक'' वापरून कांटोरचे
सुव्यवस्थापन तत्त्व सिद्ध केले; या स्वयंसिद्धकाच्या स्वीकारार्हतेवर मोठा वाद
झाला. 1910 ते 1913 या काळात रसेल आणि व्हाइटहेड यांच्या \emph{Principia
  Mathematica}चे तीन खंड प्रसिद्ध झाले. गणिताला तार्किक आधार देण्याचा
फ्रेगेचा कार्यक्रम त्यांनी पुढे नेला. मात्र रसेल आणि व्हाइटहेड यांना
निखळ तार्किक म्हणून समर्थनीय वाटणे कठीण अशी दोन तत्त्वे स्वीकारावी
लागली: अनंततेचे स्वयंसिद्धक आणि ``न्यूनीकरणीयतेचे'' स्वयंसिद्धक. 1900
च्या दशकात पॉँकारेने गणितातील ``स्वसमावेशक व्याख्यां''च्या वापरावर
टीका केली. 1910 च्या दशकात ब्रौवरने विमध्य सिद्धांत
($!A \lor \lnot !A$) न वापरता संपूर्ण गणिताची ``अंतःप्रज्ञावादी''
पायावर पुनर्मांडणी करण्याचा प्रस्ताव मांडायला सुरुवात केली.

खरोखरच विलक्षण काळ!{} सर्व गणित तर्कशास्त्रात रूपांतरित करण्याच्या
कार्यक्रमाला आता ``तर्कमीमांसावाद'' म्हणतात. वर सांगितलेल्या अडचणींमुळे
तो अयशस्वी ठरला, असे सामान्यतः मानले जाते. मानसिक अंतःप्रज्ञावादी
रचनांवर गणित उभारण्याच्या कार्यक्रमाला ``अंतःप्रज्ञावाद'' म्हणतात;
दैनंदिन गणितावर तो अति कठोर निर्बंध घालतो, असे मानले जाते. शतकाच्या
वळणावर, सर्वकाळातील अत्यंत प्रभावशाली गणितज्ञांपैकी एक असलेला डेव्हिड
हिल्बर्ट हा कांटोर आणि डेडेकिंड यांनी आणलेल्या नव्या, अमूर्त पद्धतींचा
ठाम समर्थक होता: ``कांटोरने आपल्यासाठी निर्माण केलेल्या स्वर्गातून
आपल्याला कोणीही हाकलू शकणार नाही.'' त्याच वेळी, या नव्या पद्धतींवरील
पायाभूत टीकेचीही त्याला जाणीव होती (विचित्र म्हणजे, आता या पद्धतींना
``अभिजात'' म्हणतात). दोन्ही गोष्टी साधण्याचा त्याने असा मार्ग सुचवला:
\begin{enumerate}
\item अभिजात पद्धती आकारिक स्वयंसिद्धके आणि नियम यांच्या साहाय्याने
  दर्शवा; गणितीय प्रश्न स्वयंसिद्धकीय पद्धतीतील !!{formula}s म्हणून दर्शवा.
\item या आकारिक निगमन-पद्धती सुसंगत आहेत हे सिद्ध करण्यासाठी सुरक्षित,
  ``सांतवादी'' पद्धती वापरा.
\end{enumerate}

पहिले उद्दिष्ट साधण्याच्या दिशेने हिल्बर्टच्या कार्याने मोठी मजल मारली.
1899 मध्ये त्याने आपल्या प्रसिद्ध \emph{भूमितीची पायाभरणी} या पुस्तकात
भूमितीसाठी हे काम केले होते. पुढील वर्षांत, हिल्बर्टने भूमितीसाठी जे
केले ते गणिताच्या इतर शाखांत करण्यासाठी तो, त्याचे अनेक विद्यार्थी आणि
सहकारी यांनी काम केले. हिल्बर्टने स्वतः अंकगणित आणि विश्लेषण यांसाठी
स्वयंसिद्धक-प्रणाली दिल्या. झर्मेलोने संचसिद्धान्ताची स्वयंसिद्धकीय मांडणी
दिली; फ्रेंकेल, स्कोलेम, फोन न्यूमन आणि इतरांनी ती पुढे विकसित केली.
1920 च्या दशकाच्या मध्यापर्यंत ``सर्व'' गणिताची स्वयंसिद्धकीय मांडणी
म्हणवणारे दोन दृष्टिकोन होते: रसेल आणि व्हाइटहेड यांचे \emph{Principia
  mathematica}, आणि पुढे झर्मेलो--फ्रेंकेल संच उपपत्ती म्हणून ओळखली
जाणारी मांडणी.

1921 मध्ये हिल्बर्टने या पद्धती सुसंगत असल्याचे सिद्ध करण्याच्या उद्देशाने
संशोधन कार्यक्रम सुरू केला. यात त्याला अनेक विद्यार्थ्यांची, विशेषतः
बेर्नाइस, आकरमान आणि पुढे गेन्ट्सेन यांची मदत झाली. या उद्दिष्टामागील
मूळ कल्पना अशी होती: गणितात विसंगतीची !!{derivation} शक्य आहे का, हा
प्रश्न चिन्हांच्या संभाव्य क्रमांबद्दलच्या संयुक्तिक प्रश्नात रूपांतरित
करायचा. म्हणजे, अंकगणित, विश्लेषण किंवा संचसिद्धान्त यांच्या
स्वयंसिद्धक-प्रणालीतील स्वयंसिद्धकांपासून, उदाहरणार्थ, $!A \land
\lnot !A$ ची योग्य !!{derivation} ठरणाऱ्या वाक्यांचे संभाव्य क्रम
तपासायचे. असा चिन्हक्रम अशक्य असल्याची सिद्धता ही स्वतः गणितीय सिद्धता
असल्यामुळे या स्वयंसिद्धकीय पद्धतींत आकारिकरीत्या मांडता आली असती.
दुसऱ्या शब्दांत, उदाहरणार्थ, अंकगणित सुसंगत आहे असे सांगणारे
$\OCon$ हे वाक्य असले असते. शिवाय, विशेषतः त्याची सिद्धता केवळ अतिशय
मर्यादित, ``सांतवादी'' साधनांनी होत असल्यास, ते वाक्य संबंधित पद्धतींत
सिद्ध करता आले पाहिजे होते.

विकसित स्वयंसिद्धक-प्रणाली प्रत्येक गणितीय प्रश्न सोडवतील, हे दुसरे
उद्दिष्ट दोन प्रकारे नेमके करता येते. पहिला प्रकार असा: गणिताच्या
स्वयंसिद्धक-प्रणालीच्या भाषेतील कोणतेही वाक्य~$!A$ घेतल्यास,
स्वयंसिद्धकांपासून $!A$ किंवा $\lnot !A$ सिद्ध करता येते. असे असल्यास,
स्वयंसिद्धकांच्या आधारे ज्यांची सिद्धता किंवा खंडन दोन्ही होत नाहीत अशी
वाक्ये, म्हणजे त्यांनी न सोडवलेले प्रश्न, उरणार नाहीत. हा गुणधर्म
असलेल्या स्वयंसिद्धक-प्रणालीला \emph{संपूर्ण} म्हणतात. अर्थात, दिलेल्या
एखाद्या वाक्यासाठी या दोन पर्यायांपैकी कोणता खरा आहे हे ठरवणे तरीही
कठीण असू शकते. पण तत्त्वतः ते ठरवण्याची पद्धत असली पाहिजे. प्रत्यक्षात,
हिल्बर्टने विचारात घेतलेल्या स्वयंसिद्धक आणि !!{derivation} पद्धतींसाठी
संपूर्णतेवरून अशी पद्धत अस्तित्वात आहे असे निष्पन्न होईल---जरी हिल्बर्टला
याची जाणीव नव्हती. दुसरा अर्थ अधिक प्रबळ मागणी करतो: दिलेले
वाक्य~$!A$ स्वयंसिद्धकांपासून निष्पन्न करता येते की नाही हे ठरवणारी
यांत्रिक, संगणकीय पद्धत असावी.

1931 मध्ये गोडेलने दोन ``अपूर्णतेची प्रमेये'' सिद्ध केली. त्यांवरून
हा कार्यक्रम यशस्वी होऊ शकत नाही, हे दिसले. गणितासाठी संपूर्ण अशी
स्वयंसिद्धक-प्रणाली नाही; विशेषतः, स्वयंसिद्धकांची सुसंगतता व्यक्त करणारे
वाक्य सिद्धही करता येत नाही आणि त्याचे खंडनही करता येत नाही.

यामुळे हिल्बर्टच्या मूळ कार्यक्रमाला घातक धक्का बसला. तरीही गणितात
अनेकदा घडते तसे, यातून संशोधनाच्या नव्या दिशाही उघडल्या. सर्व गणिताला
व्यापणारी एकच आकारिक पद्धत नसल्यास, अधिक मर्यादित पद्धती विकसित करून
त्यांत काय सिद्ध करता येते याचा अभ्यास करणे अर्थपूर्ण ठरते. या पद्धतींची
सुसंगतता सिद्ध करण्यासाठी कमी निर्बंध असलेल्या सिद्धता-पद्धती विकसित
करणे, आणि त्यांची सुसंगतता सिद्ध करणे किती कठीण आहे हे मोजण्याचे मार्ग
शोधणेही अर्थपूर्ण ठरते. गोडेलने दाखवल्याप्रमाणे, (जवळजवळ) प्रत्येक
आकारिक पद्धतीत असे प्रश्न असतात की जे ती सोडवू शकत नाही. म्हणून
दिलेल्या आकारिक पद्धतीला न सोडवता येणारे ``महत्त्वाचे'' प्रश्न शोधणे,
आणि ते सोडवण्यासाठी पद्धत किती प्रबळ असावी लागते हे ठरवणे योग्य ठरते.
आजतागायत तर्कशास्त्रज्ञ या प्रश्नांचा सिद्धता-उपपत्ती या नव्या
गणितीय शाखेत अभ्यास करत आहेत.

\end{document}
